\section{Introduction}
\label{sec:intro}

Short-term photovoltaic power forecasting is a productive research
area. Hybrid architectures combining convolutional and recurrent
layers, often with a further decomposition or gradient-boosted stage,
appear regularly in the literature, and reported accuracy improvements
are substantial \cite{nguyenmusgens2021drives}.

We coded 27 such papers, published between 2022 and 2026, on eight
dimensions of evaluation protocol: coded conservatively, with a field
recorded as not stated unless the paper explicitly stated it, and every
judgement supported by a verbatim quotation (full methodology in
Section II). Twenty-six of
27 report no skill score against any reference forecast, all 27 report
point estimates with no measure of run-to-run variance, twenty-two do
not state whether their split preserves temporal order, and seventeen
do not state whether night hours are included. These are not incidental
omissions: each is a choice that changes the number a paper reports,
and this paper measures by how much.

\subsection{What This Paper Does}
\label{sec:intro-what-this-paper-does}

We fix a leakage-controlled evaluation protocol, hold it constant, and
vary individual evaluation choices while keeping the model, the data,
and the training procedure identical: five architectures, three module
technologies, three horizons, two feature regimes and five seeds, 450
runs each on a validation and a held-out test year, 900 total. The
models are instruments for measuring protocol sensitivity, not
contributions in themselves.

Four results follow. The reference forecast accounts for 70 percent of
apparent skill at a six-hour horizon and inverts skill-versus-horizon
from monotonic to non-monotonic. Including night hours deflates
normalised error by approximately 34 percent, closed-form in the
daylight fraction. Fitting a residual correction stage on the split it
is evaluated on turns a component that costs 0.034 in skill into one
that appears to contribute 0.334. Architecture contributes an order of
magnitude less: the largest difference we measure between any two
architectures is 0.024 in skill score, against a 0.51-0.72
lagged-to-oracle gap.

All findings hold in direction on a held-out year evaluated once, after
every modelling and evaluation choice was frozen.

\subsection{Contributions}
\label{sec:intro-contributions}

1. An open, leakage-controlled benchmark harness, with explicit
   alignment assertions, documented data exclusions, and 900
   machine-readable run records; a verification performed on
   2026-08-08 regenerated all generated outputs from a clean state and
   reproduced them byte-identically.

2. A quantification of four protocol effects - reference forecast,
   night-hour inclusion, residual-stage construction, and weather
   information regime - measured with the model and data held constant,
   exceeding architectural differences by one to two orders of
   magnitude.

3. A component-wise decomposition of the convolutional-recurrent hybrid
   with residual correction, with paired significance testing and seed
   variance, finding no component that justifies its cost on this data.

4. A coded survey of evaluation practice across 27 papers, with
   verbatim supporting evidence, establishing that the choices measured
   here are predominantly unreported in the literature they affect.

\subsection{What This Paper Does Not Claim}
\label{sec:intro-not-claim}

The evidence base is three co-located arrays at a single site, sharing
one weather station and therefore not independent, evaluated over one
held-out year at horizons of one to six hours with deterministic
forecasts; we do not claim that any architecture is generally inferior,
that these effects transfer unchanged to other sites or climates, or
that the surveyed papers are incorrect - only that their reported
numbers cannot be compared without the protocol information most do not
provide. Rigorous evaluation practice exists in solar forecasting
\cite{yang2020verification, yang2019ropes} and is followed by one paper
in our sample of 27 \cite{mayer2022physmlhybrid}; the gap here is
non-adoption, not absent standards.
